% GENERATED FILE. Edit canonical lessons, then run npm run generate:books.
% canonical-source-hash: fnv1a64:ea5af49318fea64f
% canonical-lessons: ZH-C200-binguan, ZH-C200-yuding, ZH-C200-danrenjian, ZH-C200-ruzhu, ZH-C200-tuifang

\chapter{Booking a Room}
\label{ch:zh-booking-a-room}
\hlchaptermodality{200}

\begin{chapteropening}
\textbf{By the end of this chapter:} \emph{I can book a hotel room, ask for a single room, and say when I will check in and check out.}

Complete the last lesson of chapter 200: I can book a hotel room, ask for a single room, and say when I will check in and check out.
\end{chapteropening}

\section[bīnguǎn]{bīnguǎn — hotel}

\label{lesson:ZH-C200-binguan}



\begin{quote}
Before the new one: say the Mandarin for tight, then the Mandarin for to pick out.
\end{quote}

\subsection*{You'll want to know: bīnguǎn}
\emph{bīnguǎn} — \textquotedblleft{}hotel\textquotedblright{}.

\begin{quote}
\emph{wǒmen de bīnguǎn hěn gānjìng} — Our hotel is very clean.
\end{quote}

Its characters come later; for now it is a word for the ear.

\begin{cousinweb}
\emph{bīnguǎn} is built from \emph{bīn} 'a guest' and \emph{guǎn} 'a lodging', both inherited from Middle Chinese, \emph{pjin} and \emph{kwanH} (Guangyun, in Baxter's transcription), which glosses the first 'to welcome' and the second as a lodging house.
\end{cousinweb}

\begin{grammarlens}[title={What you've built}]
One more word for talking about plans.
\end{grammarlens}

\subsection*{Your turn}
\begin{itemize}
\raggedright
  \item \emph{Say it:} \emph{bīnguǎn}
  \item \emph{Say it:} \emph{bīnguǎn}, once more
  \item \emph{Say it:} \emph{wǒmen de bīnguǎn hěn gānjìng} — Our hotel is very clean.
  \item \emph{Recall:} say the Mandarin for tight, then the Mandarin for to pick out, then say \emph{bīnguǎn} again
\end{itemize}

\begin{tcolorbox}[breakable,colback=teal!4,colframe=teal!35!black,title={Before you move on}]
What does bīnguǎn mean? (\textquotedblleft{}Hotel\textquotedblright{}.) Say it once more.
\end{tcolorbox}
\section[yùdìng]{yùdìng — to book ahead}

\label{lesson:ZH-C200-yuding}



\begin{quote}
Before the new one: say the Mandarin for to pick out, then the Mandarin for hotel.
\end{quote}

\subsection*{You'll want to know: yùdìng}
\emph{yùdìng} — \textquotedblleft{}to book ahead\textquotedblright{}.

\begin{quote}
\emph{wǒ xiǎng yùdìng yí ge fángjiān} — I'd like to book a room.
\end{quote}

Its characters come later; for now it is a word for the ear.

\begin{cousinweb}
\emph{yùdìng} is built from \emph{yù} 'beforehand' and \emph{dìng} 'to order', both inherited from Middle Chinese, \emph{yoH} and \emph{dengX} (Guangyun, in Baxter's transcription), which glosses the first 'before'.
\end{cousinweb}

\begin{grammarlens}[title={What you've built}]
One more word for talking about plans.
\end{grammarlens}

\subsection*{Your turn}
\begin{itemize}
\raggedright
  \item \emph{Say it:} \emph{yùdìng}
  \item \emph{Say it:} \emph{yùdìng}, once more
  \item \emph{Say it:} \emph{wǒ xiǎng yùdìng yí ge fángjiān} — I'd like to book a room.
  \item \emph{Recall:} say the Mandarin for to pick out, then the Mandarin for hotel, then say \emph{yùdìng} again
\end{itemize}

\begin{tcolorbox}[breakable,colback=teal!4,colframe=teal!35!black,title={Before you move on}]
What does yùdìng mean? (\textquotedblleft{}To book ahead\textquotedblright{}.) Say it once more.
\end{tcolorbox}
\section[dānrénjiān]{dānrénjiān — single room}

\label{lesson:ZH-C200-danrenjian}



\begin{quote}
Before the new one: say the Mandarin for hotel, then the Mandarin for to book ahead.
\end{quote}

\subsection*{You'll want to know: dānrénjiān}
\emph{dānrénjiān} — \textquotedblleft{}single room\textquotedblright{}.

\begin{quote}
\emph{wǒ yào yí ge dānrénjiān} — I'd like a single room.
\end{quote}

Its characters come later; for now it is a word for the ear.

\begin{cousinweb}
\emph{dānrénjiān} is built from \emph{dān} 'single', \zh{人} 'person' and \zh{间} 'a room'. The first two are inherited from Middle Chinese, \emph{tan} and \emph{nyin} (Guangyun, in Baxter's transcription).
\end{cousinweb}

\begin{grammarlens}[title={What you've built}]
One more word for talking about plans.
\end{grammarlens}

\subsection*{Your turn}
\begin{itemize}
\raggedright
  \item \emph{Say it:} \emph{dānrénjiān}
  \item \emph{Say it:} \emph{dānrénjiān}, once more
  \item \emph{Say it:} \emph{wǒ yào yí ge dānrénjiān} — I'd like a single room.
  \item \emph{Recall:} say the Mandarin for hotel, then the Mandarin for to book ahead, then say \emph{dānrénjiān} again
\end{itemize}

\begin{tcolorbox}[breakable,colback=teal!4,colframe=teal!35!black,title={Before you move on}]
What does dānrénjiān mean? (\textquotedblleft{}Single room\textquotedblright{}.) Say it once more.
\end{tcolorbox}
\section[rùzhù]{rùzhù — to check in (at a hotel)}

\label{lesson:ZH-C200-ruzhu}



\begin{quote}
Before the new one: say the Mandarin for to book ahead, then the Mandarin for single room.
\end{quote}

\subsection*{You'll want to know: rùzhù}
\emph{rùzhù} — \textquotedblleft{}to check in (at a hotel)\textquotedblright{}.

\begin{quote}
\emph{wǒmen xiàwǔ liǎng diǎn kěyǐ rùzhù} — We can check in at two in the afternoon.
\end{quote}

Its characters come later; for now it is a word for the ear.

\begin{cousinweb}
\emph{rùzhù} is built from \emph{rù} 'to enter' and \emph{zhù} 'to stay', both inherited from Middle Chinese, \emph{nyip} and \emph{drjuH} (Guangyun, in Baxter's transcription), which glosses the first 'within' and the second 'to stop'.
\end{cousinweb}

\begin{grammarlens}[title={What you've built}]
One more word for talking about plans.
\end{grammarlens}

\subsection*{Your turn}
\begin{itemize}
\raggedright
  \item \emph{Say it:} \emph{rùzhù}
  \item \emph{Say it:} \emph{rùzhù}, once more
  \item \emph{Say it:} \emph{wǒmen xiàwǔ liǎng diǎn kěyǐ rùzhù} — We can check in at two in the afternoon.
  \item \emph{Recall:} say the Mandarin for to book ahead, then the Mandarin for single room, then say \emph{rùzhù} again
\end{itemize}

\begin{tcolorbox}[breakable,colback=teal!4,colframe=teal!35!black,title={Before you move on}]
What does rùzhù mean? (\textquotedblleft{}To check in (at a hotel)\textquotedblright{}.) Say it once more.
\end{tcolorbox}
\section[tuìfáng]{tuìfáng — to check out (of a hotel)}

\label{lesson:ZH-C200-tuifang}



\begin{quote}
Before the new one: say the Mandarin for single room, then the Mandarin for to check in.
\end{quote}

\subsection*{You'll want to know: tuìfáng}
\emph{tuìfáng} — \textquotedblleft{}to check out (of a hotel)\textquotedblright{}.

\begin{quote}
\emph{míngtiān jǐ diǎn tuìfáng} — What time is check-out tomorrow?
\end{quote}

Its characters come later; for now it is a word for the ear.

\begin{cousinweb}
\emph{tuìfáng} is built from \emph{tuì} 'to give back' and \zh{房} 'room', both inherited from Middle Chinese, \emph{thwojH} and \emph{bjang} (Guangyun, in Baxter's transcription), which glosses the first 'to draw back'.
\end{cousinweb}

\begin{grammarlens}[title={What you've built}]
One more word for talking about plans.
\end{grammarlens}

\subsection*{Your turn}
\begin{itemize}
\raggedright
  \item \emph{Say it:} \emph{tuìfáng}
  \item \emph{Say it:} \emph{tuìfáng}, once more
  \item \emph{Say it:} \emph{míngtiān jǐ diǎn tuìfáng} — What time is check-out tomorrow?
  \item \emph{Recall:} say the Mandarin for single room, then the Mandarin for to check in, then say \emph{tuìfáng} again
\end{itemize}

\begin{tcolorbox}[breakable,colback=teal!4,colframe=teal!35!black,title={Before you move on}]
What does tuìfáng mean? (\textquotedblleft{}To check out (of a hotel)\textquotedblright{}.) Say it once more.
\end{tcolorbox}
